Let $A\subseteq\mathbb{N}$ be infinite and let $f:A\to\{-1,1\}$ be non-constant.
Say that $(A,f)$ has the “zero-sum” property if there exists a finite nonempty $S\subset A$ such that
\[\sum_{n\in S}\frac{f(n)}{n}=0.\tag{$\ast$}\]

Question 1: If $A$ is an infinite arithmetic progression, must $(A,f)$ have the zero-sum property for every non-constant $f$?

Question 2: If $A$ has positive density, must $(A,f)$ have the zero-sum property for every non-constant $f$? (Statement says: no.)

Question 3: If $A=\{k^2:k\ge 2\}$ (squares excluding $1$), must $(A,f)$ have the zero-sum property?
